% The 5-cell toy column used throughout §1.6 tutorial.
\begin{tikzpicture}[
  font=\footnotesize,
  cell/.style={rectangle, draw=black, line width=0.5pt,
               minimum width=22mm, minimum height=10mm, fill=teal!4,
               anchor=north},
  surfline/.style={line width=1.2pt, color=coral},
  basalarrow/.style={-{Stealth[length=2.4mm]}, line width=0.9pt, color=forest},
  every node/.style={align=center}
]
% Surface line above cell 0
\draw[surfline] (-22mm,0) -- (22mm,0);
\node[anchor=south west, color=coral, font=\sffamily\small]
  at (-22mm, 1pt) {surface  ($T_s$ from Newton)};
% The 5 cells
\node[cell] (c0) at (0,0)         {cell 0\\$z=0.05$\,m};
\node[cell, below=0mm of c0]  (c1)        {cell 1\\$z=0.15$\,m};
\node[cell, below=0mm of c1]  (c2)        {cell 2\\$z=0.25$\,m};
\node[cell, below=0mm of c2]  (c3)        {cell 3\\$z=0.35$\,m};
\node[cell, below=0mm of c3]  (c4)        {cell 4\\$z=0.45$\,m};
% Basal flux arrow rising from below
\draw[basalarrow] (0,-58mm) -- (0,-52mm);
\node[anchor=north, color=forest, font=\sffamily\small]
  at (0,-58mm) {$Q_b = 21$\,mW\,m$^{-2}$};
% Properties box on the right
\node[anchor=north west, font=\footnotesize, color=char]
  at (32mm, 0) {
    \begin{tabular}{@{}l@{\;}l@{}}
    $\Delta z$        & $0.10$\,m \\
    $K$               & $10$\,mW\,m$^{-1}$\,K$^{-1}$ \\
    $\rho c_p$        & $9\!\times\!10^{5}$\,J\,m$^{-3}$\,K$^{-1}$ \\
    $\Delta t$        & $3600$\,s ($1$\,h) \\
    $\alpha$          & $4\!\times\!10^{-3}$ \\
    \end{tabular}
  };
\end{tikzpicture}
